\begin{tabular}{r p{0.24\linewidth} r p{0.52\linewidth}}
\toprule
\# & Persona name(s) & Conv. & Case brief (first sentence) \\
\midrule
1 & Luisa & 89 & Luisa ist genervt von den Streitereien ihrer Eltern. \\
2 & Michael Schlosser & 70 & Herr Schlosser m\"ochte gerne wieder eine Arbeit finden. \\
3 & Elke & 61 & Elke macht sich Sorgen um ihren Sohn. \\
4 & Jessica Bergmann & 53 & Jessica ist sich unsicher. \\
5 & Marlene Grafe & 43 & Sie sucht eine besser bezahlte Stelle im Hauswirtschaftbereich oder in der Gastronomie in Teilzeit (30 Stunden). \\
6 & Marco & 29 & Marco hat Schwierigkeiten, sich selbst zu organisieren. \\
7 & Eddie, Jakob, Lina, Susi, Tanja M\"uller & 28 & Lina befindet sich derzeit im Entzug. \\
8 & Frau Wolke & 25 & Frau Wolkes Sohn Jan ist 14 Jahre alt und wohnt zuhause bei ihr. \\
9 & Eddie & 19 & Eddie macht sich Sorgen um seinen Sohn. \\
10 & Frau Schuster & 13 & Frau Schusters Sohn Max ist 8 Jahre alt geht zur Schule und ist oft anschlie\ss{}end im Hort. \\
11 & Michaela & 9 & Verbesserung der Kommunikation innerhalb der Familie \\
12 & Regina & 6 & Persona arbeitet 20 Jahre im gleichen Job und kann nicht mehr die T\"atigkeit ausf\"uhren, hat k\"orperliche Beschwerden, wenn sie nur an die Arbeit denkt; m\"ochte\textbackslash{}ldots \\
13 & Lars & 6 & Lars befindet sich derzeit im Entzug. \\
14 & Boris & 5 & Familienberatung suchen, um eine neue Perspektive auf famili\"are Konflikte zu gewinnen. \\
15 & Lara & 4 & Lara ``funktioniert zwar nach au\ss{}en hin gut'', f\"uhlt sich aber innerlich unter Druck und \"uberfordert. \\
16 & Mira & 4 & Mira f\"uhlt sich zunehmend emotional ersch\"opft und \"uberfordert im Studium. \\
17 & Leyla & 3 & Leyla steht kurz nach dem Abitur und ist unsicher, welchen beruflichen Weg sie einschlagen m\"ochte. \\
18 & Annika & 3 & Annika kommt oft erst sp\"at abends nach Hause in ihre 2-Zimmer-Wohnung in Berlin. \\
19 & Michael M\"uller & 2 & Bew\"altigung der emotionalen Auswirkungen der Scheidung \\
20 & Ute Meier & 2 & Unterst\"utzung bei der Organisation eines stabilen Alltags nach der Scheidung. \\
21 & Robert 12 & 2 & Probleme des Soiihnen \\
22 & Charlie & 1 & ``Ich wei\ss{} nicht mehr weiter … es macht einfach keinen Sinn mehr.'' Charlie hat ihr Mediendesign-Studium vor wenigen Monaten abgebrochen. \\
23 & Elias & 1 & Elias hat das Gef\"uhl, dass sein aktuelles Studium keine sichere Perspektive ist -- er f\"uhlt sich unter Druck sich gegen die Konkurrenz durchzusetzen und hat\textbackslash{}ldots \\
24 & Jonas & 1 & Jonas steht kurz vor dem Studienabschluss -- er muss noch ein Seminar absolvieren und seine Bachelorarbeit schreiben. \\
25 & Marie & 1 & Trauer \\
26 & Lars & 1 & Unterst\"utzung bei der \"Uberwindung der Gl\"ucksspielsucht \\
27 & Tobias Meier & 1 & Tobias sucht Unterst\"utzung bei seiner Alkoholsucht. \\
\bottomrule
\end{tabular}
